%!TEX root=../../note
\subsection{Functions}
\subsubsection{Define functions in Racket}
How to define 
\begin{equation*}
    f(x) = x^2 + 3x + 4
\end{equation*}
in the Racket. 

\begin{minted}{rust}
    #lang htdp/bsl
    (define (f x) (+ (sqr x) (* 3 x) 4))
\end{minted}

\begin{remark}
    When you apply the function, all the parameter get replace at the first time. 
\end{remark}

\begin{definition}
    [Predicates]
    Functions that produce Bool values.
    By convention, the names of predicates end in a "?".  
\end{definition}
\begin{example}
    \begin{minted}{rust}
        (define (hot? t) (> t 30))
    \end{minted}
\end{example}

\begin{definition}[Identifiers]
Function names and parameters are identifiers and must follow certain rules.

\begin{itemize}
    \item Identifiers may contain letters, numbers, and special characters such as
    \texttt{-}, \texttt{\_}, \texttt{.}, \texttt{?}, and \texttt{=}.

    \item Identifiers cannot contain spaces, brackets of any kind, or quotation marks.

    \item Identifiers must contain at least one non-numeric character.
\end{itemize}

Parameter order matters. 
\begin{minted}{rust}
    (define (g x y) (- x y))
    (define (g y x) (- x y))
\end{minted}

Function definition are always in simplest form. 

Constant definition is simplest when the left is in simplest form.

"And" and "or" are short circuits. They are not functions. 

\end{definition}
